
\section{Pareto-Optimal Static Configurations}
\label{sec:pareto}

We quantify the inherent tradeoff resulting from the use of static
resource configuration in an environment where the application load is
unpredictable.  We specifically focus on latency-sensitive
applications with service-level agreement. 

\subsection{Energy-efficiency}


%% put the key last to have correct numbering

\begin{figure*}

\centering
 \subfloat[Linux 3.16.1]{
  \label{fig:pareto-eff:linux}
   \includegraphics{figs/blank-3-2.eps}}
 \hspace{.01in}
 \subfloat[\ix]{
  \label{fig:pareto-eff:ix}
  \includegraphics{figs/blank-3-2.eps}}

\caption{Pareto Efficiency measuring energy-efficiency as a function of throughput.}
\label{fig:pareto-eff}
\end{figure*}



Fig.~\ref{fig:pareto-eff} quantifies the tradeoff in terms of
energy-efficiency for two different operating systems, Linux 3.16.1
and \ix.  In these graphs, each black curve represents the energy
spent by the CPU, as measured by Intel RAPL~\cite{intel:rapl}, as a
function of the application load, for a given static configuration.  A
static configuration is determined by a number of execution threads
(from one to 16) and the CPU frequency of the socket (from 1.2 Ghz to
2.4Ghz, as well as Turbo mode).  The application is memcache, the load
generator mutilate~\cite{url:mutilate} and the workload is
USR~\cite{Atikoglu:2012:WAL}.  We only report results where the
application meets a service-level agreement of 500\microsecond
measured at the 99th percentile at the client. See
\S\ref{sec:eval:setup} for details on the experimental methodology and
setup.

In Fig.~\ref{fig:pareto-eff:linux} and Fig.~\ref{fig:pareto-eff:ix},
the red curve corresponds to the Pareto-optimal energy consumed for
any given application load.  The concept of Pareto efficiency (or
optimality) is broadly used in finance, life-sciences, and engeenering
to determine the optimal allocation of resources that maximizes a
particular objective function~\cite{wikipedia-pareto}.  The approach
is first useful to rule out the combinations that are suboptimal under
particular circumstances.  Fig.~\ref{fig:pareto-eff:linux} and
Fig.~\ref{fig:pareto-eff:ix} only shows the static configurations that
contribute at least one data point to to the Pareto-optimal curve.
There are 17 such static configurations for Linux and 21 for \ix.  To
put this in perspective, the cross-product of 16 possible number of
threads and 14 different clock rates leads to 224 distinct
configurations.  And of course, in the absence of any static
configuration tradeoff, a single configuration would suffice.

A detailed comparison of the Pareto-optimal curves of
Fig.~\ref{fig:pareto-eff:linux} and Fig.~\ref{fig:pareto-eff:ix}
further shows that the two operating systems have different
energy-efficiency characteristics.  First, Linux relies on interrupts
to save energy a low loads, resulting in a large dynamic spread for
each curve; in contrast, \ix relies on polling to deliver load
latency, resutling in flatter curves.  Consequently, there is also a
tradeoff between the two operating systems.  Indeed, Linux is more
energy-efficent than \ix with loads of less than XXX requests per
second.  Second, certain Linux configurations draw noticeably more
energy than \ix.  This is simply due to the fact that memcache
scalability bottlenecks at the application level limit the number of
useful threads to 9 on \ix but to 16 on Linux.  Third, the maximal
throughput of the both Linux and \ix noticeably exceed our previoiusly
reported results in \~cite{ix-osdi}; those are due primarily to recent
changes made to memcache to improve its horizontal scalability.  And
finally, at maximal load, \ix outperforms Linux by 2.0$\times$ while
drawing only 71\% of the energy.


\subsection{Server Consolidation}

Fig.~\ref{fig:pareto-consolidation} shows the similar Pareto
efficiency in the server consolidation scenario, which combines a
CPU-intensive background job with the same latency-sensitive memcache
workload.

